\section{Theory: the calculus and the algorithms that decide it}\label{sec:theory}

This section specifies the type theory that the kernel implements and the algorithms that decide it.
It is a specification of an implementation, not the presentation of a new type theory: every rule below corresponds to code in the kernel (\cref{sec:impl}) or in the independent checker (\cref{sec:indep}), and where the implementation does something other than the textbook rule we say so; the cases are collected in \cref{sec:meta}.
The calculus belongs to the Calculus-of-Inductive-Constructions family \citep{coquand88cc,coquand90ind,paulin93}.
Its nearest relatives are the kernel calculi of Lean and Rocq \citep{carneiro19,rocq,sozeau20}; the respects in which it is smaller are listed in \cref{sec:limits}.

\subsection{Conventions and configurations}\label{sec:conventions}

Typing is written $\Gamma \vdash t : A$ and definitional equality (\emph{conversion}) $\Gamma \vdash t \conv u$; the global environment $\Sg$ is fixed and left implicit.
A \emph{type} is a term $A$ with $\Gamma \vdash A : \Type_i$ for some $i$.
Substitution $\subst{t}{x}{s}$ is capture-avoiding, $\vs{a}$ abbreviates a sequence $a_1,\dots,a_n$, and $f\,\vs a$ is iterated application.
The implementation represents terms in a locally nameless style, with bound variables as de Bruijn indices and context variables as de Bruijn levels \citep{debruijn72,mcbride04,chargueraud12}; that is an implementation matter (\cref{sec:impl}) and plays no role in the rules.

The kernel has two configurations of the bottom universe, chosen when a world is created.
\begin{description}
\item[Set-like (\code{prop\_elim} off).] $\Type_0$ is impredicative and proofs are terms like any others.
  An inductive family in $\Type_0$ may not have a field in a larger universe; this is the condition under which an impredicative bottom universe admits strong elimination without paradox \citep{coquand86girard,hurkens95} (rule \textsc{A4} in \cref{sec:ind}).
\item[Prop (\code{prop\_elim} on).] $\Type_0$ is written $\Prop$.
  It is impredicative and \emph{proof-irrelevant} (\cref{sec:prop}).
  A family in $\Prop$ may have large fields, and one that is not a syntactic subsingleton can be eliminated only into $\Prop$ (\cref{sec:ind}).
  Quotient types are available (\cref{sec:quot}).
\end{description}
Every library and the whole problem bank use the second configuration, together with the option \code{parameters}, which makes the kernel distinguish a family's parameters from its indices (\cref{sec:ind}).
Rules and conditions that belong to one configuration only are marked.

\subsection{Syntax}\label{sec:syntax}

Terms, which serve also as types and as motives, are generated by
\begin{align*}
t,u,A,B,M,s \;::=\;\;
  & x \;\mid\; c \;\mid\; \Type_k \;\mid\; \Pit{x}{A}{B} \;\mid\; \lam{x}{A}{t} \;\mid\; t\,u \;\mid\; \letin{x}{A}{t}{u}\\
  \mid\;\; & \fixx{d}{f}{A}{t} \;\mid\; \matchx{s}{M}{b_1,\dots,b_n} \;\mid\; \meta{m}
\end{align*}
where $k,d\in\mathbb N$ and $c$ ranges over global names.
The kernel has eleven node kinds; ten appear above, and the eleventh, a universe-parameter node, is scaffolding for universe polymorphism that is not used (evaluation refuses it).
The term $\fixx{d}{f}{A}{t}$ is a local structural fixpoint that binds $f$ in $t$ and decreases on the $d$-th argument ($0$-based) of its type $A$ (\cref{sec:guard}).
The term $\matchx{s}{M}{b_1,\dots,b_n}$ is case analysis on $s$ with motive $M$; it has one branch per constructor of the family of $s$, in declaration order, and each branch is a curried function of \emph{all} the fields of its constructor, parameters included (\cref{sec:ind}).
The term $\meta{m}$ is a metavariable; it occurs during elaboration only, and no term admitted to the global environment contains one (\cref{sec:trust}).
Constructors are ordinary constants.
There are no primitive projections and no primitive recursors: recursion is $\sy{fix}$ together with $\sy{match}$.

A context is $\Gamma ::= \cdot \mid \Gamma, x{:}A$.
The global environment $\Sg$ is a list of declarations, each with a closed type:
\begin{itemize}
\item an \emph{axiom} $c : A$;
\item a \emph{definition} $c : A \mathrel{:=} v$, theorems included;
  every definition is transparent to conversion, and there is no opacity attribute;
\item an \emph{inductive family} with its constructors (\cref{sec:ind});
\item a \emph{recursive function} $f : A \mathrel{:=} t$ with a decreasing argument (\cref{sec:guard});
\item the five \emph{quotient constants} (\cref{sec:quot}).
\end{itemize}

\subsection{Typing}\label{sec:typing}

\Cref{fig:rules-core} gives the rules for the variable, universe, dependent-product and application fragment, and \cref{fig:rules-data} those for local definitions, fixpoints and case analysis.

\begin{figure}[t]
\centering
\begin{mathpar}
\inferrule*[right=\pf{Sort}]{k < 2^{30}}{\Gamma \vdash \Type_k : \Type_{k+1}}
\and
\inferrule*[right=\pf{Var}]{(x{:}A)\in\Gamma}{\Gamma \vdash x : A}
\and
\inferrule*[right=\pf{Const}]{(c : A)\in\Sg}{\Gamma \vdash c : A}
\\
\inferrule*[right=\pf{Pi}]{\Gamma \vdash A : \Type_i \\ \Gamma, x{:}A \vdash B : \Type_j}{\Gamma \vdash \Pit{x}{A}{B} : \Type_{\imax(i,j)}}
\and
\inferrule*[right=\pf{Lam}]{\Gamma \vdash A : \Type_i \\ \Gamma, x{:}A \vdash t : B}{\Gamma \vdash \lam{x}{A}{t} : \Pit{x}{A}{B}}
\\
\inferrule*[right=\pf{App}]{\Gamma \vdash f : \Pit{x}{A}{B} \\ \Gamma \vdash a : A}{\Gamma \vdash f\,a : \subst{B}{x}{a}}
\and
\inferrule*[right=\pf{Conv}]{\Gamma \vdash t : A \\ \Gamma \vdash A \conv A' \\ \Gamma \vdash A' : \Type_i}{\Gamma \vdash t : A'}
\end{mathpar}
\[
\imax(i,0) = 0,\qquad \imax(i,j) = \max(i,j)\ \ (j>0).
\]
\caption{Typing rules, part one.
  The side condition of \textsc{Sort} keeps the successor of a level from wrapping around in the 32-bit counter that stores it.
  \textsc{Pi} is the same in both configurations: the bottom universe is impredicative.}
\label{fig:rules-core}
\end{figure}

\paragraph{Universes.}
The hierarchy is $\Type_0 : \Type_1 : \Type_2 : \cdots$, with levels that are numerals.
There is \emph{no cumulativity}: a term of type $\Type_i$ does not have type $\Type_j$ for $j\neq i$, and a coercion between levels is an explicit construction.
There is also \emph{no universe polymorphism}: the calculus has no level variables, no level constraints and no level expressions.
A development that needs a definition at several levels instantiates a construction function once per level vector, and the kernel checks each instance (\cref{sec:impl}).
This is a deliberate restriction, taken to keep the universe fragment of the kernel trivially decidable and the independent checker small; its cost is discussed in \cref{sec:limits}.
The side condition $k<2^{30}$ in \textsc{Sort} is a guard against a specific failure: if levels wrapped around in 32-bit arithmetic, $\Type_{2^{32}-1}$ would be assigned the type $\Type_0$, and a universe that contains itself admits Girard's paradox \citep{coquand86girard,hurkens95}.

\paragraph{Dependent products.}
The rule \textsc{Pi} places $\Pit{x}{A}{B}$ in $\Type_{\imax(i,j)}$.
When the codomain is a proposition ($j=0$) the product is a proposition whatever the level of the domain: this is impredicativity of the bottom universe \citep{coquand88cc}, and it is what allows $\forall (A:\Type_5).\,P\,A$ to be a proposition.
The kernel computes the level of a product telescope $\Pi(x_1{:}A_1)\cdots(x_n{:}A_n).B$ from the inside out, starting from the level of $B$ and applying $\imax$ with the level of each $A_j$.

\paragraph{Application.}
The kernel types a whole application spine at once.
The type of the head is exposed as a product telescope, by unfolding definitions at the head on syntax and, only if that fails, by evaluating and reading back; each argument is checked against its domain with the earlier arguments substituted; and the result type is obtained by substitution.
Types are not normalised on the way.
This is the rule \textsc{App} iterated, with \textsc{Conv} applied where a domain must be matched.

\paragraph{Bidirectional operation.}
The rules are declarative; the kernel decides them bidirectionally.
Inference synthesises a type for a term; checking verifies a term against a given type, and falls back on inference followed by conversion.
Checking is where the two non-trivial elaboration steps happen: a $\sy{match}$ checked against an expected type gets its motive by abstracting the scrutinee and its indices from that type (\cref{sec:ind}), and a $\lambda$ takes its expected domain from the expected product.
Neither step is trusted: the term that results is checked in full, with the motive that was built.

\begin{figure}[t]
\centering
\begin{mathpar}
\inferrule*[right=\pf{Let}]{\Gamma \vdash A : \Type_i \\ \Gamma \vdash v : A \\ \Gamma \vdash \subst{u}{x}{v} : C}{\Gamma \vdash \letin{x}{A}{v}{u} : C}
\\
\inferrule*[right=\pf{Fix}]{\Gamma \vdash A : \Type_i \\ \Gamma, f{:}A \vdash t : A \\ \sy{guard}_d(f,t,A)}{\Gamma \vdash \fixx{d}{f}{A}{t} : A}
\\
\inferrule*[right=\pf{Match}]{
  \Gamma \vdash s : I\,\vs{a} \qquad
  \Gamma \vdash M : \Pi(\vs{y}{:}\vs{Y})(z{:}Z).\,S \qquad |\vs y| = |\vs a| \qquad S \text{ a sort} \\
  \Gamma \vdash M\,\vs{a}\,s : S \qquad
  \sy{Elim}(I,S) \\
  \Gamma \vdash \sy{Br}(I,c,\vs{a},M) : \Type_{j_c} \qquad
  \Gamma \vdash b_c : \sy{Br}(I,c,\vs{a},M)\ \ (1\le c\le n)}
 {\Gamma \vdash \matchx{s}{M}{b_1,\dots,b_n} : M\,\vs{a}\,s}
\end{mathpar}
\caption{Typing rules, part two.
  $\sy{guard}_d$ is the structural-recursion condition of \cref{sec:guard}.
  $\sy{Br}$, the type of a branch, and $\sy{Elim}$, the elimination restriction, are defined in \cref{sec:ind}.
  $I$ is the family of the scrutinee, with constructors $C_1,\dots,C_n$.}
\label{fig:rules-data}
\end{figure}

\paragraph{Local definitions.}
\textsc{Let} types the body with the value substituted: the context never contains a definitional entry, and $\sy{let}$ is $\zeta$-reduced rather than abstracted.
The cost is that a body that uses its let-bound variable many times is re-typed against the value each time; the benefit is that contexts consist of assumptions only.

\subsection{Reduction and conversion}\label{sec:conv}

Reduction is the union of the following rules, closed under all term formers.
\begin{align*}
(\beta)\quad & (\lam{x}{A}{t})\,u \;\red\; \subst{t}{x}{u}\\
(\zeta)\quad & \letin{x}{A}{v}{u} \;\red\; \subst{u}{x}{v}\\
(\delta)\quad & c \;\red\; v \qquad\text{if } c : A \mathrel{:=} v \in \Sg\\
(\iota)\quad & \matchx{C_c\,\vs{v}}{M}{b_1,\dots,b_n} \;\red\; b_c\,\vs{v}\\
(\varphi)\quad & (\fixx{d}{f}{A}{t})\;a_0 \cdots a_d\,\vs{a}' \;\red\; \subst{t}{f}{\fixx{d}{f}{A}{t}}\;a_0\cdots a_d\,\vs{a}'
   \qquad\text{if } a_d = C\,\vs{w}\text{ for a constructor }C\\
(\kappa)\quad & \Quot.\sy{lift}\,A\,r\,B\,f\,h\,(\Quot.\sy{mk}\,A'\,r'\,a) \;\red\; f\,a
\end{align*}
In $(\iota)$ the branch is applied to \emph{all} fields $\vs v$ of the constructor, parameters included, which is why a branch is a curried function of all of them.
The body of a recursive function is a $\lambda$-abstraction, so $(\delta)$ yields a value and does not diverge: a recursive call inside a branch is not evaluated until its branch is selected.
The rule $(\kappa)$ is the only reduction rule that is not $\beta$, $\delta$, $\zeta$, $\iota$ or the unfolding of a fixpoint (\cref{sec:quot}).

\emph{Conversion} $\Gamma \vdash t \conv u$ is the least congruence that contains $\red^*$ and is closed under function $\eta$,
\[
\lam{x}{A}{t\,x} \;\conv\; t\qquad (x\notin\FV(t)),
\]
and, in the Prop configuration, under definitional proof irrelevance (\cref{sec:prop}).
There is no $\eta$ for records or for the unit type, and none for a $\sy{fix}$ beyond the function $\eta$ above.

\paragraph{Deciding conversion.}
Conversion is decided by normalisation by evaluation \citep{berger91,gregoire02,abel18}, in an environment-based, strict style.
Terms are evaluated, in an environment of values, into a semantic domain $\mathcal V$:
\begin{align*}
v \;::=\;\; & \langle\lam{x}{A}{t},\rho\rangle \mid \langle\Pit{x}{A}{B},\rho\rangle \mid \Type_k \mid C\,\vs v \mid \langle\fixx{d}{f}{A}{t},\rho\rangle \mid h\,\vs\phi\\
h \;::=\;\; & \text{a context level} \mid \text{a global constant} \mid \meta{m}
\qquad
\phi \;::=\; \mathsf{ap}(v) \mid \mathsf{case}(M,\vs b,\rho)
\end{align*}
Closures pair a term with its environment $\rho$; a \emph{neutral} term $h\,\vs\phi$ is a head applied to a spine of frames (an application, or a stuck case analysis that keeps its motive, branches and environment).
The evaluator satisfies
\begin{align*}
\ev(\rho, t\,u) &= \sy{apply}(\ev(\rho,t),\ev(\rho,u)),\\
\sy{apply}(\langle\lam{x}{A}{t},\rho\rangle, v) &= \ev((\rho,x{\mapsto}v),t),\\
\ev(\rho,\letin{x}{A}{v}{u}) &= \ev((\rho,x{\mapsto}\ev(\rho,v)),u),\\
\ev(\rho,\matchx{s}{M}{\vs b}) &= \sy{elim}(\ev(\rho,s),\langle M,\vs b,\rho\rangle),\\
\sy{elim}(C_c\,\vs v,\langle M,\vs b,\rho\rangle) &= \sy{apply}^*(\ev(\rho,b_c),\vs v),\\
\ev(\rho,c) &= \ev(\cdot,v)\ \text{ if } c:A\mathrel{:=}v\in\Sg,\qquad \ev(\rho,c)=c\ \text{ otherwise,}
\end{align*}
together with the rules that saturating a constructor produces a value $C\,\vs v$, that a $\sy{case}$ frame is added to a neutral scrutinee, that $\Quot.\sy{lift}$ fires $(\kappa)$ when its sixth argument is a saturated $\Quot.\sy{mk}$, and that an application of a fixpoint value unfolds one layer.
$\Gamma\vdash t\conv u$ is decided by running $\mathsf{conv}_n(\ev(\rho_\Gamma,t),\ev(\rho_\Gamma,u))$, where $n=|\Gamma|$ and $\rho_\Gamma$ maps each context variable to the neutral term at its level:
\begin{itemize}
\item two closures, or two product closures (after their domains), are compared by instantiating both bodies with the neutral term at the fresh level $n$ and comparing at $n+1$;
\item a closure against a neutral term or a fixpoint is $\eta$-expanded: the other side is applied to the fresh neutral term and the results are compared;
\item two constructor values must have the same constructor and pairwise convertible fields;
\item two neutral terms must have the same head and spines of equal length whose frames are pairwise convertible; two stuck case frames must have the same motive term and branches that are convertible in their captured environments.
\end{itemize}
The comparison of stuck case frames and of fixpoints in captured environments is sound but incomplete, in the usual way: a pair that is convertible but was built in differently shaped environments may be rejected.
When a spine comparison fails at a pair of frames that are not both applications, the failure is \emph{uncertain}, and the comparison is repeated on the read-back normal forms of the two neutral terms; this recovers the cases that matter in practice, at the price of normalising.
The environment-based evaluator is strict in its arguments; since the calculus is intended to be strongly normalising (\cref{sec:meta}), the strategy affects cost and not the relation.
A fixpoint is unfolded one layer when it is applied, whatever the argument is, while the specification $(\varphi)$ unfolds it only at a constructor; the independent checker follows the specification.
The difference does not show in the artefacts we re-check: the independent checker accepts every global of every library and every reference proof of the bank (\cref{sec:indep}), and a proof that a tactic builds from a fixpoint wraps it in a barrier constant that keeps it folded (finding F23 in \cref{sec:testing}).
In general, whether the eager unfolding can make the kernel accept a conversion that the specification does not is part of the gap discussed in \cref{sec:meta}.
